Table~\ref{tab:alpha} and Fig.~\ref{fig:alpha} (left) are the gain slice of the grid, Fig.~\ref{fig:alpha} (right) the period slice; safety slices are in Table~\ref{tab:grid}. Cells are means over 5 seeds.

\begin{table}[t]\centering\caption{Gain sweep at $\Delta t=0.05$\,s, cells ``time to goal [s] / min.\ clearance [m]''.}\label{tab:alpha}
\footnotesize\setlength{\tabcolsep}{3pt}\begin{tabular}{lcccc}
\hline
$\alpha$ & CT-CBF (DI) & CT-CBF (Uni) & DT-CBF (DI) & DT-CBF (Uni) \\
\hline
0.5 & 10.93 / 0.95 & 12.84 / 0.40 & -- / 0.80 & 12.76 / 0.38 \\
1 & 8.80 / 0.56 & 12.17 / 0.23 & 8.62 / 0.43 & 12.09 / 0.21 \\
2 & 7.80 / 0.20 & 11.86 / 0.15 & 7.54 / 0.13 & 11.78 / 0.14 \\
5 & 7.30 / 0.03 & 11.88 / 0.14 & 6.93 / 0.01 & 11.80 / 0.13 \\
10 & 7.05 / 0.01 & 11.92 / 0.14 & -- / 0.00 & 11.85 / 0.13 \\
\hline
\end{tabular}
\end{table}
\begin{figure*}[t]\centering\includegraphics[width=0.46\textwidth]{sweep_alpha.pdf}\hfill\includegraphics[width=0.46\textwidth]{sweep_dt.pdf}\caption{Effect of $\alpha$ at $\Delta t=0.05$\,s (left) and of the control period $\Delta t$ at $\alpha=2$ (right); mean $\pm$ sample std over seeds, violation rate at barrier level.}\label{fig:alpha}\end{figure*}

\textbf{Gain.} Large $\alpha$ makes DT-CBF unsafe on the double integrator (violation rate 0.222 at $\alpha=5$, 0.618 at $\alpha=10$, Table~\ref{tab:grid}) while CT-CBF stays at zero on this slice. On the slice the depth of the DT-CBF violations is small (mean worst penetration 0.003\,m at $\alpha=10$), so the binary violation rate overstates the physical severity there compared with the braking heuristic (0.184\,m at the default). \textbf{Period.} On the slice at $\alpha=2$ CT-CBF is insensitive to $\Delta t$ up to 0.2\,s, while DT-CBF's violation rate goes from 0.004 at $\Delta t\le0.02$ to 0.070 at 0.2; the steps up to $\Delta t=0.1$ are within the seed spread shown in Fig.~\ref{fig:alpha}, and the last step (0.018 to 0.070) exceeds it.

\textbf{Between samples or at samples?} Table~\ref{tab:samp} separates episodes with a barrier-level violation at any 0.005\,s check from episodes with one at the control sample instants, with the same $10^{-9}$\,m tolerance. \emph{Double integrator, DT-CBF.} The at-sample rate is 0.000 in every cell of Table~\ref{tab:samp} except $\alpha=10,\Delta t=0.2$ (0.002, one episode in 500), including the saturated cells where the any-time rate is 1.000. All other violations of this filter, the deep ones at $\alpha\Delta t\ge1$ included, therefore occur between samples: the robot is outside the obstacles at the samples and the zero-order-hold path dips inside in between. With $\gamma=1$ the condition only requires the successor to be on the boundary and does not limit the approach speed, which may be why the dips are deep there (not isolated). For $\alpha\Delta t<1$ the any-time rate is up to 0.618 ($\alpha=10,\Delta t=0.05$) and the mean worst penetration up to 0.016\,m ($\alpha=2,\Delta t=0.2$, Table~\ref{tab:grid}); it is 0.000\,m at three decimals only for $\Delta t\le0.02$. At $\Delta t=0.005$\,s, equal to the plant step, there is no between-sample check and the rate is 0.000 for $\alpha=2,5,10$. The reported rates therefore depend on the 0.005\,s check resolution: the condition does not give safety between samples, but almost no violation is seen at the samples. The single at-sample episode was not investigated. \emph{Double integrator, CT-CBF.} Among the gains of Table~\ref{tab:samp} it violates at samples only at $\alpha=10$ (0.018 of 0.270 at $\Delta t=0.1$, 0.184 of 0.386 at $\Delta t=0.2$); at $\alpha=5,\Delta t=0.2$ its violations (0.020) are between samples only. The cell $\alpha=0.5,\Delta t=0.2$, not in Table~\ref{tab:samp}, also has at-sample violations in the run registry. \emph{Unicycle.} For DT-CBF the two rates coincide in every cell but one (0.734 against 0.732), and the violations remain at $\Delta t=0.005$\,s (0.034, 0.826 and 0.992 for $\alpha=2,5,10$): they are not inter-sample dips, the look-ahead point is inside the inflated obstacle at the samples themselves. Their mean worst penetration grows with $\Delta t$ (0.001, 0.002, 0.004, 0.009, 0.023\,m at $\alpha=5$, Table~\ref{tab:grid}). This is consistent with the mismatch between the Euler successor of the look-ahead point used by the filter and the motion actually executed, which may be the cause; an exact-successor variant was not run.

\begin{table*}[t]\centering\caption{Violation rate at any 0.005\,s check / at control sample instants only (barrier level, tolerance $10^{-9}$\,m), means over 5 seeds. $\Delta t=0.005$\,s equals the plant step. Unicycle CT-CBF (not shown) is at 0.000 / 0.000 for $\Delta t\le0.1$ including 0.005, and both rates equal the Table~\ref{tab:grid} rate at $\Delta t=0.2$.}\label{tab:samp}
\scriptsize\setlength{\tabcolsep}{4pt}\begin{tabular}{lllcccccc}
\hline
Task & Filter & $\alpha$ & $\Delta t$=0.005 & $\Delta t$=0.01 & $\Delta t$=0.02 & $\Delta t$=0.05 & $\Delta t$=0.1 & $\Delta t$=0.2 \\
\hline
DI & CT-CBF & 2 & 0.000 / 0.000 & 0.000 / 0.000 & 0.000 / 0.000 & 0.000 / 0.000 & 0.000 / 0.000 & 0.000 / 0.000 \\
DI & CT-CBF & 5 & 0.000 / 0.000 & 0.000 / 0.000 & 0.000 / 0.000 & 0.000 / 0.000 & 0.000 / 0.000 & 0.020 / 0.000 \\
DI & CT-CBF & 10 & 0.000 / 0.000 & 0.000 / 0.000 & 0.000 / 0.000 & 0.000 / 0.000 & 0.270 / 0.018 & 0.386 / 0.184 \\
\hline
DI & DT-CBF & 2 & 0.000 / 0.000 & 0.004 / 0.000 & 0.004 / 0.000 & 0.008 / 0.000 & 0.018 / 0.000 & 0.070 / 0.000 \\
DI & DT-CBF & 5 & 0.000 / 0.000 & 0.078 / 0.000 & 0.108 / 0.000 & 0.222 / 0.000 & 0.402 / 0.000 & 1.000 / 0.000 \\
DI & DT-CBF & 10 & 0.000 / 0.000 & 0.358 / 0.000 & 0.434 / 0.000 & 0.618 / 0.000 & 1.000 / 0.000 & 1.000 / 0.002 \\
\hline
Uni & DT-CBF & 2 & 0.034 / 0.034 & 0.092 / 0.092 & 0.184 / 0.184 & 0.464 / 0.464 & 0.734 / 0.732 & 0.970 / 0.970 \\
Uni & DT-CBF & 5 & 0.826 / 0.826 & 0.900 / 0.900 & 0.960 / 0.960 & 0.988 / 0.988 & 1.000 / 1.000 & 1.000 / 1.000 \\
Uni & DT-CBF & 10 & 0.992 / 0.992 & 0.996 / 0.996 & 1.000 / 1.000 & 1.000 / 1.000 & 1.000 / 1.000 & 1.000 / 1.000 \\
\hline
\end{tabular}
\end{table*}

\textbf{Threshold.} The braking heuristic's violations fall with $d_{th}$ (Table~\ref{tab:dth}) and reach zero at $d_{th}=2$ on the double integrator, but there the success rate is 0.970, and since a timeout counts as 30\,s this raises its time to goal (9.11\,s). \textbf{Obstacle selection.} Choosing the nearest obstacle instead of the most critical one raises the double-integrator violation rate of CT-CBF from 0.00 to 0.81 and leaves that of DT-CBF at 0.01 (Table~\ref{tab:sel}). For CT-CBF this choice changes safety far more than $\Delta t$ does at the default gain.

\begin{table}[t]\centering\caption{Braking threshold sweep, cells ``violation rate / worst penetration [m] / success rate / time to goal [s]'' (barrier level).}\label{tab:dth}
\footnotesize\setlength{\tabcolsep}{3pt}\begin{tabular}{lcc}
\hline
$d_{th}$ & Braking (DI) & Braking (Uni) \\
\hline
0.25 & 1.00 / 0.657 / 1.000 / 6.50 & 0.00 / 0.000 / 1.000 / 11.97 \\
0.5 & 0.87 / 0.531 / 1.000 / 7.03 & 0.00 / 0.000 / 1.000 / 12.10 \\
1 & 0.11 / 0.184 / 1.000 / 7.48 & 0.00 / 0.000 / 1.000 / 12.57 \\
2 & 0.00 / 0.000 / 0.970 / 9.11 & 0.00 / 0.000 / 1.000 / 13.43 \\
\hline
\end{tabular}
\end{table}
\begin{table}[t]\centering\caption{Obstacle selection (double integrator, $\alpha=2$, $\Delta t=0.05$), cells ``violation rate / worst penetration [m] / time to goal [s]''.}\label{tab:sel}
\footnotesize\begin{tabular}{llcc}
\hline
System & Task & sel=nearest & sel=critical \\
\hline
CT-CBF & DI & 0.81 / 0.666 / 6.63 & 0.00 / 0.000 / 7.80 \\
DT-CBF & DI & 0.01 / 0.000 / 7.33 & 0.01 / 0.000 / 7.54 \\
\hline
\end{tabular}
\end{table}
